\section{Unterer Bauraum}

Im folgenden werden alle Berechnungen, die für die Wellen im unteren Bauraum benötigt waren aufgeführt. Für jegliche Berechnungen wurde Kysssoft verwendet. Es werden dennoch die Ergebnisse und die dazugehörigen Konstruktionsansätze hier kurz benannt. Zur Nachvollziehung der gesamten Berechnungen sind die hinzugefügten  Kysssoft-Dateien heranzuziehen. 
\\
\subsection{Zahnräderauslegung}
Zahnräder 


\subsection{Wellenauslegung}
Bei der Wellenauslegung wurde berücksichtigt, dass die Lagerlebensdauer, die Dauerfestigkeit der Wellen, die statischen Sicherheit und die Sicherheiten für die Welle-Nabe-Verbindungen eingehalten werden. 

\subsubsection{Lagerlebensdauer}
Um die Lagerlebensdauer nach ISO 281 zu berechnen muss zuerst die Länge der Lebensdauer festgelegt werden. Dabei wird zwischen der Aufwendigkeit des Austschens der Lager und der Anforderung, dass die Lager längstmöglich halten, abgewogen (QUELLE). Hierbei ist zu berücksichtigen, dass eine längere Lagerlebensdauer zu teureren und größeren Lagern führt (QUELLE). Aufgrund jener Abwägung wurde die Lagerlebensdauer auf zehn Jahre, ca. 30.000 Stunden festgelegt. 




